% TAFSD-614: Attendance Machines API (ZKTeco ADMS push), internal reference. Built with Tectonic:
%   tectonic -X compile docs/api-docs/attendance-machines.tex --outdir docs/api-docs/pdf
% (see docs/api-docs/README.md). Keep in step with src/modules/api-docs/content/attendance-machines.ts;
% src/modules/api-docs/api-docs.spec.ts checks the device list against ZK_DEVICES.
\documentclass[11pt]{article}
\usepackage{tafs-api-doc}
\tafsheader{TAFS API \,--\, Attendance Machines}

\title{\textbf{TAFS API Documentation} \\[4pt]
\large Attendance Machines \\ ZKTeco biometric devices (ADMS push / iclock)}
\author{The American Foundation School (TAFS) \\ TAFS IT and Software Department}
\date{Version 1.0 (internal) \\ 27 November 2026}

\begin{document}
\begin{titlepage}
\vspace*{\fill}
{\let\newpage\relax\maketitle}
\thispagestyle{empty}
\vspace*{\fill}
\internalonly
\end{titlepage}

\tableofcontents
\newpage

% ---------------------------------------------------------------------------
\section{Overview}
Each campus has ZKTeco biometric devices (NYU series, and SenseFace models on Push v3 firmware). They speak ZKTeco's
\textbf{ADMS push protocol}: the device opens every connection, TAFS never polls it. The device registers, fetches its
settings, then posts each punch as it happens and polls about once a minute for commands.
\begin{enumerate}[nosep]
  \item The raw request is stored in \texttt{zk\_push\_logs} (body capped at 64,000 characters). It shows under
        \textbf{ZK Device Logs $\rightarrow$ Raw logs}.
  \item An \texttt{ATTLOG} push is parsed line by line into \texttt{zk\_attendance\_scans}, one row per punch.
  \item The pair (device serial, PIN) is looked up in \texttt{device\_user\_mappings} to find the employee or student.
  \item That person's day is rebuilt: IN/OUT directions, \texttt{attendance\_staff\_daily} or
        \texttt{attendance\_student\_daily}.
  \item A \textbf{live} punch notifies the employee (Staff App) or the student's parents. Backfill never notifies.
\end{enumerate}

\subsection{Known devices}
From \texttt{ZK\_DEVICES} in \texttt{zk-device-health.util.ts}:
\begin{center}
\begin{tabularx}{\textwidth}{@{}l P{0.27\textwidth} l l L@{}}
\toprule
\textbf{Serial} & \textbf{Name} & \textbf{Campus} & \textbf{Sat.} & \textbf{Students it may map} \\
\midrule
\texttt{NYU7261205221} & Campus 2 Device 1 Secondary & GEJ & No & GEJ, Secondary \\
\texttt{NYU7261205141} & Campus 2 Device 2 Senior Cambridge & GEJ & No & GEJ, Senior Cambridge \\
\texttt{NYU7261205172} & TAFSAL & GEJ & Yes & GEJ, O / A Levels \\
\texttt{NYU7261205142} & Campus 3 Device 1 Junior Cambridge & GEJ & No & GEJ, Junior Cambridge \\
\texttt{NYU7261205128} & Campus 3 Device 2 Pre-Primary & GEJ & No & GEJ, Early Years \\
\texttt{NYU7251000240} & Johar Faculty & GEJ & Yes & Any GEJ student \\
\texttt{NYU7261000023} & NNN Faculty & NNZ & Yes & Any NNZ student \\
\texttt{NYU7261205040} & GKF Faculty & KNF & Yes & Any KNF student \\
\bottomrule
\end{tabularx}
\end{center}
Names are the on-box labels. ``Campus 2'' and ``Campus 3'' devices are physically at the Johar (GEJ) campus, so the
campus column comes from scan history, not the label.

% ---------------------------------------------------------------------------
\section{Setting Up a Device}
Every endpoint lives under the API's global prefix, so the device must reach
\url{https://api.tafs.cloud/api/v1/iclock/...}. The \texttt{/api/v1} prefix has applied to these routes since
30 June 2026.
\begin{enumerate}[nosep]
  \item On the device, open the cloud server (ADMS) settings under the communication menu. Set the server address to
        the TAFSync API host and turn on the domain-name option.
  \item Set the device clock to Pakistan time (UTC+05:00). Punch times are stored exactly as the device reports them,
        and the server sends a \texttt{Date} header on \texttt{registry}, \texttt{cdata} and \texttt{regstable} for the
        device to sync from.
  \item Add the serial to \texttt{ZK\_DEVICES} (\texttt{zk-device-health.util.ts}), and to
        \texttt{DEVICE\_SEGMENT\_RULES} or \texttt{DEVICE\_CAMPUS\_ONLY\_RULES} (\texttt{zk-device-segment-rules.ts}) if
        students will use it. Keep the webapp's \texttt{zk-devices.ts} in step. Until then the device works, but it gets
        no health tile and no student eligibility check.
  \item Enrol people on the device. A student's PIN must be their CC. An employee's PIN can be anything, then map it
        (Section~\ref{sec:mapping}).
  \item Check \textbf{ZK Device Logs $\rightarrow$ Raw logs}: a \texttt{registry} row, then \texttt{cdata} rows. Unknown
        PINs appear under unmapped PINs, with the name the device reports.
\end{enumerate}
\tafsnote{Menu wording differs between firmware versions, and some models do not accept a path in the server address.
Confirm on the model being installed that its requests arrive at \texttt{/api/v1/iclock/...} (Raw logs shows the first
\texttt{registry} row) before handing it over.}

% ---------------------------------------------------------------------------
\section{Authentication}
There is \textbf{none}. Device firmware cannot send credentials, so the \texttt{/iclock/*} routes are open. A device is
identified only by the \texttt{SN} query parameter (\texttt{sn} is also accepted; a request with neither is logged as
\texttt{unknown}).
\begin{itemize}[nosep]
  \item A punch from an unknown serial, or an unmapped PIN, is stored but credited to nobody.
  \item A forged request carrying a real serial and a mapped PIN would be credited. Nothing in the protocol prevents it;
        keep the URL out of public material.
  \item Bodies are read as plain text whatever the \texttt{Content-Type}, up to \textbf{2 MB} (\texttt{main.ts}). Larger
        bodies are refused by the server with HTTP 413.
  \item Maintenance mode does not block these routes (\texttt{app-maintenance.guard.ts}).
\end{itemize}

% ---------------------------------------------------------------------------
\section{Endpoints}
Base path \texttt{https://api.tafs.cloud/api/v1/iclock}. Every route answers HTTP 200 with a \texttt{text/plain} body.
Failures inside TAFS are logged, never returned: the device always gets its \texttt{OK} so it does not retry forever.
\begin{center}
\begin{tabularx}{\textwidth}{@{}l P{0.35\textwidth} L P{0.17\textwidth}@{}}
\toprule
\textbf{Method} & \textbf{Path} & \textbf{Sent by the device when} & \textbf{Response} \\
\midrule
POST & \texttt{/registry?SN=...} & It starts up (ADMS v2 / Push v3) & \texttt{OK} + \texttt{Date} \\
GET & \texttt{/cdata?SN=...\&options=all} & It asks for its settings & Settings block \\
POST & \texttt{/regstable?SN=...} & Push v3: it declares its tables & \texttt{OK} + \texttt{Date} \\
POST & \texttt{/cdata?SN=...\&table=ATTLOG} & It sends punches, user changes, photos & \texttt{OK} \\
GET & \texttt{/getrequest?SN=...} & It polls for commands (about every 60 s) & \texttt{OK}, or a command \\
POST & \texttt{/devicecmd?SN=...} & It reports a command's result & \texttt{OK} \\
POST & \texttt{/querydata?SN=...} & Some Push v3 firmware answers a data query here & \texttt{OK} (raw only) \\
\bottomrule
\end{tabularx}
\end{center}

\subsection{POST /registry}
Body: comma-separated \texttt{key=value} pairs describing the hardware and firmware. Stored raw. The reply must be
exactly \texttt{OK}; anything else is not a valid ADMS v2 command and makes the device retry indefinitely.
\begin{lstlisting}
POST /api/v1/iclock/registry?SN=NYU7251000240

DeviceType=acc,~DeviceName=NYU,FirmVer=...,PushVersion=...
\end{lstlisting}
\begin{lstlisting}
HTTP/1.1 200 OK
Content-Type: text/plain
Date: Fri, 27 Nov 2026 03:00:00 GMT

OK
\end{lstlisting}

\subsection{GET /cdata (settings)}
The reply is fixed. Each line matters to some firmware:
\begin{lstlisting}
OK
Stamp=0
ATTLOGStamp=0
OPERLOGStamp=0
ATTPHOTOStamp=0
Delay=60
TransFlag=TransData AttLog OpLog AttPhoto
TimeOut=10
Realtime=1
Encrypt=None
PushOptionsFlag=1
\end{lstlisting}
\begin{itemize}[nosep]
  \item \texttt{Stamp=0} and the per-table stamps tell the device to send everything it holds from the beginning.
        Without them, devices with \texttt{LogIDFunOn=1} (SenseFace 2A) never enter the push phase. The cost is a full
        backfill after every settings fetch, which the duplicate rules (Section~\ref{sec:processing}) make harmless.
  \item \texttt{Realtime=1} sends each punch as it happens; \texttt{Delay=60} is the batch interval otherwise.
  \item \texttt{PushOptionsFlag=1} must be echoed for Push v3 firmware, or it loops back to \texttt{registry}.
\end{itemize}

\subsection{POST /cdata (punches and other tables)}
The \texttt{table} query parameter says what the body holds. TAFS answers \texttt{OK} straight away and processes the
body afterwards, because a backfill can carry thousands of lines and a slow reply makes the device resend. At most two
pushes are processed at once.
\begin{center}
\begin{tabularx}{\textwidth}{@{}l L@{}}
\toprule
\textbf{table} & \textbf{What TAFS does} \\
\midrule
\texttt{ATTLOG} & Each line becomes a punch (format below). \\
\texttt{OPERLOG} & Lines like \texttt{USER PIN=16\textbackslash tName=Ali Khan\textbackslash t...} save the device's name
  for that PIN in \texttt{zk\_pin\_name\_hints}. Used only to suggest names on the unmapped-PIN screen. \\
anything else & \texttt{USERINFO}, \texttt{ATTPHOTO}, \texttt{BIOPHOTO}, \ldots: stored raw in \texttt{zk\_push\_logs},
  nothing more. \\
\bottomrule
\end{tabularx}
\end{center}
An \texttt{ATTLOG} line is tab-separated, in one of two layouts. TAFS tells them apart by where the date-time sits:
\begin{lstlisting}
# LogIDFunOn=0 (ADMS v1):  PIN  DateTime  Status  Verify  WorkCode
8088	2026-11-27 07:42:15	0	1	0

# LogIDFunOn=1 (ADMS v2):  LogID  PIN  DateTime  Status  Verify  WorkCode
10452	8088	2026-11-27 07:42:15	0	1	0
\end{lstlisting}
\begin{center}
\begin{tabularx}{\textwidth}{@{}l l L@{}}
\toprule
\textbf{Field} & \textbf{Stored as} & \textbf{Notes} \\
\midrule
PIN & \texttt{device\_pin} & Required. With the serial, decides who punched. \\
DateTime & \texttt{scan\_time} & \texttt{YYYY-MM-DD HH:MM:SS}, device wall clock, stored as written (no time-zone conversion). \\
Status & \texttt{device\_status} & The device's in/out key. Stored, but \textbf{not used}: direction is worked out by TAFS. \\
Verify & \texttt{verify\_mode} & How the person was verified (finger, face, card, \ldots). Stored only. \\
WorkCode & \texttt{work\_code} & Stored only. \\
\bottomrule
\end{tabularx}
\end{center}
A line with no PIN or an unreadable date-time is skipped silently.

\subsection{GET /getrequest and POST /devicecmd (commands)}
TAFS can queue \textbf{one} command per device, held in memory. The next poll receives it as
\texttt{C:<id>:<command>}; otherwise the reply is \texttt{OK}. The device reports the result to \texttt{devicecmd},
which is stored with a \texttt{DEVICECMD} prefix in the raw logs. A newer command replaces one never picked up, and a
restart drops whatever is queued.
\begin{lstlisting}
GET /api/v1/iclock/getrequest?SN=NYU7251000240

C:482113:DATA QUERY ATTLOG StartTime=2026-11-27 00:00:00	EndTime=2026-11-27 23:59:59
\end{lstlisting}
\begin{lstlisting}
POST /api/v1/iclock/devicecmd?SN=NYU7251000240

ID=482113&Return=0&CMD=DATA
\end{lstlisting}
\texttt{Return=0} means success. The only command TAFS sends today is the nightly re-send request
(Section~\ref{sec:quirks}).

% ---------------------------------------------------------------------------
\section{From Serial and PIN to a Person}\label{sec:mapping}
A PIN means nothing on its own: the same number can be different people on different devices.
\texttt{device\_user\_mappings} holds one row per \textbf{(device serial, PIN)}, unique, pointing at either an employee
(\texttt{person\_type = STAFF}, \texttt{employee\_id}) or a student (\texttt{STUDENT}, \texttt{student\_cc}). Someone
enrolled on three devices needs three rows.
\begin{itemize}[nosep]
  \item An inactive mapping, or none, leaves the punch stored but credited to nobody. It shows on the unmapped-PIN
        screen.
  \item Creating, editing, deactivating or deleting a mapping re-credits \textbf{every past punch} on that serial and
        PIN, then rebuilds the affected days. It runs inline for up to 2,000 punches. Beyond that the mapping is saved,
        history is not rebuilt, and the response says to finish it with:
\begin{lstlisting}
POST /api/v1/attendance/zk-scan-resolution/resolve
{"kind":"device_pin", "device_sn":"...", "device_pin":"...", "dry_run":false}
\end{lstlisting}
  \item Posting over a deactivated mapping is refused (409) unless \texttt{is\_active: true} is sent.
  \item Mappings are managed in \textbf{ZK Device Logs} (mappings, unmapped PINs, PIN lookup) and in the biometric tabs
        of the Employee and Student directories.
\end{itemize}
Checks run when a mapping is saved. A hard block cannot be overridden; a block or warning can be acknowledged with
\texttt{acknowledge\_collisions}.
\begin{center}
\begin{tabularx}{\textwidth}{@{}l l L@{}}
\toprule
\textbf{Code} & \textbf{Level} & \textbf{Meaning} \\
\midrule
\texttt{PIN\_NOT\_EQUAL\_TO\_CC} & Hard block & A student's PIN must equal their CC. \\
\texttt{DEVICE\_SEGMENT\_MISMATCH} & Hard block & The student's campus or segment is not one this device serves. \\
\texttt{PIN\_IS\_OTHER\_STUDENT\_CC} & Block & The PIN is another student's CC; their punches would be credited here. \\
\texttt{PIN\_IS\_OTHER\_STUDENT\_GR} & Block & The PIN is another student's GR number. GR numbers and CCs share one number space. \\
\texttt{PIN\_USED\_ON\_OTHER\_DEVICE} & Warning & The same PIN is mapped to a different person on another device. \\
\bottomrule
\end{tabularx}
\end{center}

% ---------------------------------------------------------------------------
\section{How a Punch Is Processed}\label{sec:processing}
\begin{longtable}{@{}p{0.2\textwidth} p{0.74\textwidth}@{}}
\toprule
\textbf{Rule} & \textbf{Detail} \\
\midrule
\endhead
Exact duplicate & Unique key (serial, PIN, time). A punch already stored is dropped before anything else runs, so
  replaying a whole day is safe. \\
Double tap & A punch within \textbf{2 minutes} of the same person's last accepted punch, on any device, is kept with
  \texttt{is\_duplicate = true} and does not change their IN/OUT sequence. \\
Live or backfill & A punch within \textbf{10 minutes} of the current Pakistan time is live (\texttt{is\_live}). Only live
  punches notify anyone, and only live punches go to Coral9's attendance feed. \\
Staff direction & Strict alternation through the day: 1st IN, 2nd OUT, 3rd IN, \ldots \\
Student direction & Against the class's internal \textbf{intermediate time}: every punch before it is an IN (only the
  first is the arrival), the first at or after it is an OUT, later ones alternate. Classes with no intermediate time
  fall back to alternation. \\
Student lateness & Paused: \texttt{STUDENT\_LATE\_MARKING\_ENABLED = false}, so every student arrival counts as
  PRESENT. \\
Protected days & A student day on leave, or marked by hand or by the system as anything but ABSENT, is not overwritten
  by a punch. A default ABSENT mark is. \\
Left students & A punch by a student who is LEFT, EXPELLED or GRADUATED is processed normally and also flagged for the
  Not Enrolled check-ins report. \\
Parent notice & Sent for the first arrival and for departures. Repeat taps before the intermediate time are silent. \\
\bottomrule
\end{longtable}

% ---------------------------------------------------------------------------
\section{Known Quirks}\label{sec:quirks}
\begin{longtable}{@{}p{0.2\textwidth} p{0.74\textwidth}@{}}
\toprule
\textbf{Quirk} & \textbf{What to know} \\
\midrule
\endhead
Clock skew & Times are taken as the device states them. A device whose clock drifts shifts its punches by the same
  amount. The old standalone campus device ran \textbf{2 minutes behind} Johar Faculty (\texttt{NYU7251000240}),
  measured over 2,649 punches seen on both. The 2-minute double-tap window uses a strict \texttt{<}, so it does not
  catch a skew of exactly 2 minutes: de-duplicate such imports per person per day, never per timestamp. \\
\texttt{OLDDEV-XLS} & Synthetic serial reserved for back-filling the old standalone device's Excel export
  (27 Jul -- 30 Aug 2026), with times shifted \textbf{+2 minutes} to match. Each person needs a mapping on this serial;
  5 exist. As of 27 Nov 2026 no punches have been imported under it. \\
\texttt{MANUAL} & Pseudo-serial for gate-desk check-ins a staff member punches in the webapp for a student with no
  biometric enrolment. They run through the same pipeline (directions, parent notices); the serial is what tells them
  apart. Re-crediting never touches these rows. \\
Full backfill on reconnect & \texttt{Stamp=0} makes a device resend all it holds after each settings fetch. Duplicates
  are dropped by the unique key; raw bodies over 64,000 characters are truncated in \texttt{zk\_push\_logs} (the
  punches themselves are all processed). \\
Devices that never flush & On 12 Sep 2026 the API was down 05:3x--10:44 PKT. Four devices sent their queued punches on
  reconnect; Johar Faculty and GKF Faculty did not. Trial fix: every day at 23:00 PKT TAFS queues
  \texttt{DATA QUERY ATTLOG} for that day on those two devices only (\texttt{zk-device-resend-scheduler.service.ts}). \\
Push v3 firmware & Needs \texttt{PushOptionsFlag=1} echoed in the settings and a 200 on \texttt{regstable}. Either
  missing sends it back to \texttt{registry} in a loop. \\
Device health & Expected to talk Monday to Saturday, 06:00--18:00 PKT. After 15 minutes of silence the tile shows a
  warning, after 60 an alert, after 120 critical. On Saturdays only the faculty devices and TAFSAL are expected. Last
  contact includes the minute-by-minute polls, which are kept in memory only, so straight after a restart health falls
  back to the last stored push. \\
Three device lists & \texttt{ZK\_DEVICES}, the \texttt{DEVICE\_*\_RULES} maps and the webapp's \texttt{zk-devices.ts}
  must be edited together. \\
\bottomrule
\end{longtable}

% ---------------------------------------------------------------------------
\section{Error Codes}
The device endpoints have no error codes: every accepted request gets HTTP 200 and \texttt{OK}. The only failures a
device can see come from the web server itself:
\begin{center}
\begin{tabularx}{\textwidth}{@{}l L@{}}
\toprule
\textbf{HTTP} & \textbf{When} \\
\midrule
200 & Every request to a known \texttt{/iclock} route, including ones whose processing failed (see the server log and
  Raw logs). \\
404 & A route TAFS does not implement. Firmware that gets one typically falls back to \texttt{registry}. \\
413 & Body larger than 2 MB. \\
\bottomrule
\end{tabularx}
\end{center}
The staff-facing mapping routes under \path{/api/v1/attendance/zk-device-mappings} return the codes in
Section~\ref{sec:mapping} with HTTP 409, and 400 for a STAFF mapping without \texttt{employee\_id} or a STUDENT mapping
without \texttt{student\_cc}.

% ---------------------------------------------------------------------------
\section{Testing}
There is no sandbox: every device and every test talks to production.
\begin{itemize}[nosep]
  \item \textbf{Simulate scan} (ZK Device Logs, action \texttt{attendance.zk\_device\_logs\#simulate},
        \texttt{POST} \path{/api/v1/attendance/zk-device-mappings/simulate-scan} with \texttt{device\_sn}, \texttt{device\_pin},
        optional \texttt{scan\_time}) runs one punch through the real pipeline and returns the stored punch and the
        rebuilt day.
  \item To try the raw protocol, post to \texttt{/cdata} with a made-up serial. Its punches stay unmapped and change
        nobody's attendance.
\end{itemize}
\begin{lstlisting}
curl -X POST 'https://api.tafs.cloud/api/v1/iclock/cdata?SN=TEST-DEVICE-01&table=ATTLOG' \
  --data-binary $'99999\t2026-11-27 09:00:00\t0\t1\t0'
\end{lstlisting}
\tafsnote{\textbf{Simulate scan always notifies}, even for a past time. With a real serial and a mapped PIN it marks
that person's real attendance and messages them or their parents. Use a test student or an unmapped PIN.}

\end{document}
